\documentclass[11pt]{article}
\usepackage[letterpaper,margin=0.72in]{geometry}
\usepackage[T1]{fontenc}
\usepackage[utf8]{inputenc}
\usepackage{lmodern}
\usepackage{microtype}
\usepackage{xcolor}
\usepackage{graphicx}
\usepackage{booktabs}
\usepackage{longtable}
\usepackage{tabularx}
\usepackage{array}
\usepackage{enumitem}
\usepackage[most]{tcolorbox}
\usepackage{fancyhdr}
\usepackage{titlesec}
\usepackage{hyperref}
\usepackage{xurl}
\usepackage{lastpage}
\usepackage{ragged2e}
\usepackage{multicol}
\usepackage{amssymb}
\definecolor{BAOrange}{HTML}{E87521}
\definecolor{BANavy}{HTML}{17324D}
\definecolor{BASlate}{HTML}{5B6573}
\definecolor{BALight}{HTML}{F4F6F8}
\definecolor{BAGold}{HTML}{D7A62A}
\hypersetup{colorlinks=true,linkcolor=BANavy,urlcolor=BAOrange,citecolor=BANavy,pdftitle={Requirements and Traceability Templates}}
\setlength{\parindent}{0pt}
\setlength{\parskip}{4pt}
\setlist{nosep,leftmargin=*}
\renewcommand{\arraystretch}{1.18}
\newcolumntype{Y}{>{\RaggedRight\arraybackslash}X}
\newcolumntype{L}[1]{>{\RaggedRight\arraybackslash}p{#1}}
\pagestyle{fancy}
\setlength{\headheight}{14pt}
\fancyhf{}
\lhead{\textcolor{BANavy}{\small BABOK v3 Study Series}}
\rhead{\textcolor{BASlate}{\small Requirements and Traceability Templates}}
\cfoot{\small \thepage\ / \pageref{LastPage}}
\renewcommand{\headrulewidth}{0.4pt}
\titleformat{\section}{\Large\bfseries\color{BANavy}}{\thesection}{0.7em}{}
\titleformat{\subsection}{\large\bfseries\color{BAOrange}}{\thesubsection}{0.6em}{}
\titleformat{\subsubsection}{\normalsize\bfseries\color{BANavy}}{\thesubsubsection}{0.5em}{}
\newtcolorbox{keybox}[1][]{enhanced,breakable,colback=BALight,colframe=BANavy,boxrule=0.7pt,arc=1.5mm,left=2mm,right=2mm,top=1.5mm,bottom=1.5mm,#1}
\newtcolorbox{practicebox}[1][]{enhanced,breakable,colback=white,colframe=BAOrange,boxrule=0.8pt,arc=1.5mm,left=2mm,right=2mm,top=1.5mm,bottom=1.5mm,#1}
\newtcolorbox{warningbox}[1][]{enhanced,breakable,colback=white,colframe=BAGold,boxrule=0.8pt,arc=1.5mm,left=2mm,right=2mm,top=1.5mm,bottom=1.5mm,#1}
\begin{document}
\begin{titlepage}
\centering
\vspace*{1.15in}
{\Huge\bfseries\color{BANavy} Requirements and Traceability Templates\par}
\vspace{0.35in}
{\Large\color{BAOrange} Reusable practitioner artifacts aligned to BABOK v3 concepts\par}
\vspace{0.55in}
\rule{0.78\textwidth}{1.2pt}\par
\vspace{0.35in}
{\large Technical reference and study artifact based on the uploaded \emph{A Guide to the Business Analysis Body of Knowledge (BABOK Guide), Version 3}.\par}
\vfill
\begin{keybox}
\textbf{Source discipline.} BABOK terminology and organization are preserved where practical. Explanations are paraphrased for study and engineering use. Any engineering, security, Agile, UML/BPMN, or enterprise-architecture crosswalks are practitioner mappings, not normative BABOK statements.
\end{keybox}
\vspace{0.35in}
{\small Prepared October 2026\par}
\end{titlepage}
\tableofcontents
\newpage
\section{Scope and Use}
These are original, reusable templates synthesized from the source concepts of requirement classification, traceability, maintenance, prioritization, change assessment, approval, verification, validation, and requirements architecture. They are \textbf{not official IIBA forms}. Tailor fields to the initiative and do not collect data merely because a field exists.

\section{Requirement Classification Reference}
\begin{longtable}{L{0.20\textwidth} L{0.68\textwidth}}
\toprule
Business & Why the change exists: goals, objectives, and outcomes.\\
Stakeholder & What a stakeholder needs in order for business requirements to be achieved.\\
Solution -- Functional & Behavior or information capabilities the solution must provide.\\
Solution -- Non-functional & Conditions and qualities under which the solution must remain effective.\\
Transition & Temporary capabilities or conditions needed to move from current to future state.\\
\bottomrule
\end{longtable}
\section{Template A: Atomic Requirement Record}
\begin{longtable}{L{0.25\textwidth} L{0.63\textwidth}}
\toprule
Field & Entry\\
\midrule
Requirement ID & \rule{0.93\linewidth}{0.4pt}\\
Title & \rule{0.93\linewidth}{0.4pt}\\
Class & $\square$ Business \quad $\square$ Stakeholder \quad $\square$ Functional \quad $\square$ Non-functional \quad $\square$ Transition\\
Statement & \vspace{0.35in}\\
Rationale / need & \vspace{0.28in}\\
Source stakeholder(s) & \vspace{0.20in}\\
Value / outcome supported & \vspace{0.20in}\\
Acceptance / evaluation criteria & \vspace{0.36in}\\
Priority and rationale & \vspace{0.20in}\\
Dependencies / assumptions / constraints & \vspace{0.30in}\\
Backward trace & Business/stakeholder need IDs: \rule{0.55\linewidth}{0.4pt}\\
Forward trace & Design/component/test/release IDs: \rule{0.55\linewidth}{0.4pt}\\
Status / version & \rule{0.60\linewidth}{0.4pt}\\
Owner / approver & \rule{0.60\linewidth}{0.4pt}\\
\bottomrule
\end{longtable}
\section{Template B: Non-Functional Requirement Record}
\begin{longtable}{L{0.25\textwidth} L{0.63\textwidth}}
\toprule
Field & Entry\\
\midrule
NFR ID / quality category & \vspace{0.18in}\\
Context / operating condition & \vspace{0.22in}\\
Quality statement & \vspace{0.30in}\\
Measure / unit & \vspace{0.18in}\\
Target / threshold & \vspace{0.18in}\\
Measurement method / evidence & \vspace{0.25in}\\
Failure consequence / risk & \vspace{0.24in}\\
Related functional requirements & \vspace{0.18in}\\
Design constraints, if any & \vspace{0.20in}\\
Acceptance owner & \vspace{0.18in}\\
\bottomrule
\end{longtable}

\begin{practicebox}
\textbf{Security extension.} For AppSec use, common quality categories include confidentiality, integrity, authentication, authorization, auditability, privacy, resilience, dependency integrity, secure configuration, and security observability. The category is only a label; the requirement still needs a measurable condition and evidence method.
\end{practicebox}
\section{Template C: Traceability Matrix}
\small
\begin{longtable}{L{0.10\textwidth} L{0.13\textwidth} L{0.13\textwidth} L{0.14\textwidth} L{0.14\textwidth} L{0.12\textwidth} L{0.08\textwidth}}
\toprule
Req ID & Business / stakeholder source & Related reqs & Design / component & Verification evidence & Release / state & Owner\\
\midrule
REQ-001 & BR-001 / SR-003 & REQ-004 & CMP-API-01 & TC-122 / SCAN-88 & R2 / Approved & A. Owner\\
\rule{0.8in}{0.4pt} & & & & & & \\
\rule{0.8in}{0.4pt} & & & & & & \\
\bottomrule
\end{longtable}
\normalsize

\subsection{Trace Relationship Types}
\begin{itemize}
\item \textbf{Derives from / satisfies:} links lower-level requirements to higher-level needs or requirements.
\item \textbf{Depends on / enables:} records sequencing or logical dependency.
\item \textbf{Allocated to / implemented by:} connects requirements to design or solution components.
\item \textbf{Verified by / evidenced by:} connects requirements to tests, reviews, scans, inspections, or other evidence.
\item \textbf{Conflicts with / constrains:} records relationships that create trade-offs or boundaries.
\end{itemize}
\section{Template D: Requirement Change Impact Assessment}
\begin{longtable}{L{0.28\textwidth} L{0.60\textwidth}}
\toprule
Field & Entry\\
\midrule
Change request ID / date & \\
Proposed change & \vspace{0.25in}\\
Reason / triggering need & \vspace{0.20in}\\
Affected requirements & \\
Affected stakeholders & \\
Affected designs/components/interfaces & \vspace{0.20in}\\
Affected tests / acceptance evidence & \vspace{0.20in}\\
Value impact & \vspace{0.18in}\\
Risk impact & \vspace{0.18in}\\
Cost / effort / schedule impact & \vspace{0.18in}\\
Transition / migration impact & \vspace{0.18in}\\
Options considered & \vspace{0.22in}\\
Recommendation & \vspace{0.22in}\\
Decision / approver / date & \vspace{0.18in}\\
\bottomrule
\end{longtable}
\section{Template E: Prioritization Worksheet}
\begin{longtable}{L{0.22\textwidth} L{0.14\textwidth} L{0.14\textwidth} L{0.14\textwidth} L{0.14\textwidth}}
\toprule
Requirement & Value & Urgency & Risk / uncertainty & Dependency / sequencing\\
\midrule
REQ-001 & 1--5: \rule{0.25in}{0.4pt} & 1--5: \rule{0.25in}{0.4pt} & 1--5: \rule{0.25in}{0.4pt} & Notes: \rule{0.7in}{0.4pt}\\
REQ-002 & & & & \\
REQ-003 & & & & \\
\bottomrule
\end{longtable}
\begin{warningbox}
Avoid turning prioritization into an unexplained arithmetic score. The purpose is to make trade-offs explicit and support a decision. Record the rationale, dependencies, and dissent where material.
\end{warningbox}
\section{Template F: Verification and Validation Checklist}
\begin{longtable}{L{0.48\textwidth} L{0.10\textwidth} L{0.30\textwidth}}
\toprule
Check & Pass? & Notes / evidence\\
\midrule
Statement is clear enough for its intended audience and purpose. & $\square$ & \\
Terminology is consistent with the domain glossary/data dictionary. & $\square$ & \\
Requirement does not conflict with another approved requirement without an explicit trade-off. & $\square$ & \\
Acceptance/evaluation criteria are measurable or otherwise objectively assessable. & $\square$ & \\
Dependencies, assumptions, constraints, and interfaces are represented. & $\square$ & \\
Backward trace to a business/stakeholder need exists where required. & $\square$ & \\
Requirement supports the business outcome or value it claims to support. & $\square$ & \\
Requirement is feasible enough to proceed to the next decision point. & $\square$ & \\
Required stakeholders reviewed the information. & $\square$ & \\
\bottomrule
\end{longtable}
\section{Template G: Approval Record}
\begin{longtable}{L{0.26\textwidth} L{0.62\textwidth}}
\toprule
Baseline / package & \\
Scope of approval & \vspace{0.20in}\\
Decision & $\square$ Approved \quad $\square$ Approved with conditions \quad $\square$ Rejected \quad $\square$ Deferred\\
Conditions / exceptions & \vspace{0.28in}\\
Approver(s) / authority & \vspace{0.18in}\\
Date / effective version & \\
Related change / risk decisions & \vspace{0.20in}\\
\bottomrule
\end{longtable}
\section{Template H: Requirements Architecture / Allocation Matrix}
\small
\begin{longtable}{L{0.12\textwidth} L{0.18\textwidth} L{0.18\textwidth} L{0.16\textwidth} L{0.16\textwidth} L{0.10\textwidth}}
\toprule
Req ID & Parent / source & Domain / capability & Allocated component & Related process / interface & Status\\
\midrule
REQ-001 & BR-001 & Identity & API Gateway & Login / OAuth callback & Approved\\
 & & & & & \\
 & & & & & \\
\bottomrule
\end{longtable}
\normalsize
\section{Template I: AppSec Security Requirement Trace Extension}
This extension is a practitioner mapping that preserves BABOK traceability concepts while adding security-specific evidence.
\small
\begin{longtable}{L{0.10\textwidth} L{0.13\textwidth} L{0.13\textwidth} L{0.13\textwidth} L{0.13\textwidth} L{0.13\textwidth} L{0.10\textwidth}}
\toprule
Sec Req & Asset / data & Threat / abuse case & Control / design & Code / component & Evidence & Residual risk / owner\\
\midrule
SEC-001 & Customer token & Credential theft & Short-lived token + rotation & Auth service & Unit/integration + SAST/DAST review & Low / IAM owner\\
 & & & & & & \\
 & & & & & & \\
\bottomrule
\end{longtable}
\normalsize

\section*{Primary Source}
\addcontentsline{toc}{section}{Primary Source}
International Institute of Business Analysis (IIBA), \emph{A Guide to the Business Analysis Body of Knowledge (BABOK Guide)}, Version 3, 2015. This document is a study and practitioner synthesis of the user-supplied copy; it is not an official IIBA publication.
\end{document}
